% !TEX program = xelatex
% Lernplatt - by Can Siebert
% Kurs 71 (Dig 42) - Tag 16 - Aufgabenheft (Capstone)
% CSRD / EU-Taxonomie als Prozess-, Daten- und Kontrollarchitektur
%
% Self-contained xelatex source. Kein biblatex, keine externen Bilder.

\documentclass[11pt,a4paper,oneside]{article}

\usepackage{polyglossia}
\setdefaultlanguage{german}

\usepackage[a4paper,margin=2.5cm]{geometry}

\usepackage{fontspec}
\defaultfontfeatures{Ligatures=TeX}

\IfFontExistsTF{Charter}{%
  \setmainfont{Charter}%
}{%
  \IfFontExistsTF{Bitstream Charter}{%
    \setmainfont{Bitstream Charter}%
  }{%
    \setmainfont{TeX Gyre Schola}%
  }%
}

\IfFontExistsTF{Source Sans 3}{%
  \setsansfont{Source Sans 3}%
}{%
  \IfFontExistsTF{Source Sans Pro}{%
    \setsansfont{Source Sans Pro}%
  }{%
    \setsansfont{TeX Gyre Heros}%
  }%
}

\newcommand{\caveatfontfallback}{\itshape}
\IfFontExistsTF{Caveat}{%
  \newfontfamily\caveatfont{Caveat}[Scale=1.15]%
}{%
  \newcommand\caveatfont{\caveatfontfallback}%
}

\setmonofont{Latin Modern Mono}

\usepackage{microtype}
\linespread{1.15}

\usepackage{xcolor}
\definecolor{lpInk}{RGB}{20,28,38}
\definecolor{lpAccent}{RGB}{157,74,26}
\definecolor{lpAccentLight}{RGB}{246,228,212}
\definecolor{lpAmber}{RGB}{222,138,30}
\definecolor{lpAmberLight}{RGB}{255,243,217}
\definecolor{lpAmberBorder}{RGB}{196,118,18}
\definecolor{lpGray}{RGB}{96,104,114}
\definecolor{lpGrayLight}{RGB}{238,240,243}
\definecolor{lpGreen}{RGB}{34,120,72}
\definecolor{lpGreenLight}{RGB}{222,238,228}
\definecolor{lpGreenBorder}{RGB}{24,96,58}

\definecolor{lpL1}{RGB}{34,120,72}
\definecolor{lpL2}{RGB}{22,90,140}
\definecolor{lpL3}{RGB}{170,42,52}

\usepackage[most]{tcolorbox}
\tcbuselibrary{breakable,skins}

\newtcolorbox{infobox}[1][Hinweis]{%
  enhanced, breakable,
  colback=lpAccentLight,
  colframe=lpAccent,
  boxrule=0.6pt,
  arc=2pt,
  left=10pt,right=10pt,top=8pt,bottom=8pt,
  fonttitle=\sffamily\bfseries\color{white},
  coltitle=white,
  title={#1},
  attach boxed title to top left={xshift=8pt,yshift=-8pt},
  boxed title style={size=small, colback=lpAccent, colframe=lpAccent, boxrule=0.4pt, arc=1pt},
  top=14pt
}

\newtcolorbox{antwortbox}[1][Ihre Antwort]{%
  enhanced, breakable,
  colback=white,
  colframe=lpGray,
  boxrule=0.4pt,
  arc=2pt,
  left=10pt,right=10pt,top=8pt,bottom=8pt,
  fonttitle=\sffamily\bfseries\color{white},
  coltitle=white,
  title={#1},
  attach boxed title to top left={xshift=8pt,yshift=-8pt},
  boxed title style={size=small, colback=lpGray, colframe=lpGray, boxrule=0.4pt, arc=1pt},
  top=14pt
}

\usepackage{enumitem}
\setlist{itemsep=2pt, topsep=4pt, parsep=0pt}
\setlist[itemize,1]{label={\textbullet},leftmargin=1.6em}
\setlist[itemize,2]{label={\textendash},leftmargin=1.4em}
\setlist[enumerate,1]{label=\arabic*., leftmargin=1.8em}

\usepackage{tabularx}
\usepackage{array}
\usepackage{booktabs}
\usepackage{longtable}

\usepackage{fancyhdr}
\usepackage{lastpage}
\pagestyle{fancy}
\fancyhf{}
\renewcommand{\headrulewidth}{0.3pt}
\renewcommand{\footrulewidth}{0pt}
\fancyhead[L]{\footnotesize\sffamily\color{lpGray}Aufgabenheft \textperiodcentered{} Kurs 71 \textperiodcentered{} Tag 16}
\fancyhead[R]{\footnotesize\sffamily\color{lpGray}CSRD / EU-Taxonomie \textendash{} Capstone}
\fancyfoot[L]{\footnotesize\sffamily\color{lpGray}\textcopyright{} Can Siebert}
\fancyfoot[R]{\footnotesize\sffamily\color{lpGray}\thepage{} / \pageref{LastPage}}

\fancypagestyle{plain}{%
  \fancyhf{}%
  \renewcommand{\headrulewidth}{0pt}%
  \fancyfoot[C]{\footnotesize\sffamily\color{lpGray}\thepage{} / \pageref{LastPage}}%
}

\usepackage[explicit]{titlesec}

\titleformat{\section}[block]
  {\sffamily\Large\bfseries\color{lpAccent}}
  {\thesection.}{0.6em}{#1}
  [{\vspace{2pt}\color{lpAccent}\titlerule[0.6pt]}]

\titleformat{\subsection}[block]
  {\sffamily\large\bfseries\color{lpInk}}
  {\thesubsection}{0.6em}{#1}

\titlespacing*{\section}{0pt}{14pt}{8pt}
\titlespacing*{\subsection}{0pt}{10pt}{4pt}

\usepackage[hidelinks,unicode]{hyperref}

\newcommand{\akzent}[1]{{\caveatfont\color{lpAmber}#1}}
\newcommand{\fachbegriff}[1]{\textbf{#1}}

\newcommand{\niveaubadge}[2]{%
  \tcbox[on line, colback=#1, colframe=#1, coltext=white,
         boxrule=0pt, arc=2pt, left=5pt,right=5pt,top=1pt,bottom=1pt,
         nobeforeafter]{\sffamily\bfseries\footnotesize #2}%
}
\newcommand{\nivLi}{\niveaubadge{lpL1}{L1\,\textbullet\,Wissen}}
\newcommand{\nivLii}{\niveaubadge{lpL2}{L2\,\textbullet\,Anwendung}}
\newcommand{\nivLiii}{\niveaubadge{lpL3}{L3\,\textbullet\,Management}}

\newcommand{\zeit}[1]{%
  \tcbox[on line, colback=lpGrayLight, colframe=lpGrayLight, coltext=lpInk,
         boxrule=0pt, arc=2pt, left=5pt,right=5pt,top=1pt,bottom=1pt,
         nobeforeafter]{\sffamily\footnotesize\textbf{$\sim$\,#1}}%
}

\newcommand{\aufgabenkopf}[4]{%
  \par\vspace{6pt}
  \noindent{\sffamily\Large\bfseries\color{lpAccent}Aufgabe #1 \textendash{} #2}\par
  \vspace{2pt}
  \noindent{\color{lpAccent}\rule{\linewidth}{0.6pt}}\par
  \vspace{4pt}
  \noindent #3 \hfill \zeit{#4}\par
  \vspace{6pt}
}

\newcommand{\ytquery}[1]{%
  \par\vspace{4pt}
  \noindent{\footnotesize\sffamily\color{lpGray}\textbf{YouTube-Suchquery:} \texttt{#1}}\par
}

\newcommand{\antwortlinien}[1]{%
  \par\vspace{6pt}
  \foreach \x in {1,...,#1}{%
    \noindent\makebox[\linewidth]{\dotfill}\\[6pt]
  }
}
\usepackage{pgffor}

% =========================================================================
\begin{document}

\begin{titlepage}
  \pagestyle{empty}
  \vspace*{1.5cm}
  \noindent{\sffamily\footnotesize\color{lpGray}LERNPLATT \textperiodcentered{} BY CAN SIEBERT}\\[2pt]
  \noindent{\color{lpAccent}\rule{\linewidth}{1.2pt}}\\[4pt]
  \noindent{\sffamily\footnotesize\color{lpGray}AUFGABENHEFT \textperiodcentered{} MODUL 8 (CAPSTONE) \textperiodcentered{} TAG 16 VON 16 \textperiodcentered{} KURS 71 (Dig 42) \textperiodcentered{} 19.\,Juni 2026}

  \vspace{3cm}
  {\sffamily\fontsize{30}{34}\selectfont\bfseries\color{lpInk}Aufgabenheft\\Tag 16\par}

  \vspace{0.7cm}
  {\sffamily\large\color{lpGray}CSRD \& EU-Taxonomie \textendash{}\\Prozess-, Daten- und Kontroll-\\architektur in elf Aufgaben.\par}

  \vspace{1.5cm}
  {\caveatfont\LARGE\color{lpAmber}11 Aufgaben \textperiodcentered{} L1 \textperiodcentered{} L2 \textperiodcentered{} L3 \textperiodcentered{} Capstone}\par

  \vfill

  \noindent\begin{minipage}{0.55\linewidth}
    {\sffamily\footnotesize\color{lpGray}Niveau-Mix:}\\
    {\sffamily\small\color{lpInk}3\,\textsf{L1} \textperiodcentered{} 5\,\textsf{L2} \textperiodcentered{} 3\,\textsf{L3}}\\[10pt]
    {\sffamily\footnotesize\color{lpGray}Bearbeitung:}\\
    {\sffamily\small\color{lpInk}Selbstbearbeitung im Lernpfad}
  \end{minipage}\hfill
  \begin{minipage}{0.4\linewidth}
    \raggedleft
    {\sffamily\footnotesize\color{lpGray}Dozent:}\\
    {\sffamily\small\color{lpInk}Can Siebert}\\[10pt]
    {\sffamily\footnotesize\color{lpGray}Begleitend:}\\
    {\sffamily\small\color{lpInk}Lehrskript \& L\"osungsheft}
  \end{minipage}

  \vspace{0.5cm}
  \noindent{\color{lpAccent}\rule{\linewidth}{0.6pt}}
\end{titlepage}

\section*{So arbeiten Sie mit diesem Heft}
\thispagestyle{plain}

\noindent Dieses Aufgabenheft enth\"alt die elf Capstone-Aufgaben zu Tag 16 \textendash{} \emph{Regulatorik in Steuerung \"ubersetzen}: EU-Taxonomie und CSRD als Prozess-, Daten- und Kontroll\-architektur, integriert in laufende Transformationsprogramme (Prozessdigitalisierung, RPA, Green IT, Change). Die Aufgaben sind in drei Niveaus gegliedert:

\begin{itemize}
  \item \nivLi{} \textendash{} \fachbegriff{Wissen.} CSRD-/EU-Taxonomie-Logik, Controls, Evidence Pack, KPI-Dictionary. 5\,--\,10 Minuten.
  \item \nivLii{} \textendash{} \fachbegriff{Anwendung.} Reportingprozess, Controls, KPI-Steckbriefe, Anti-Gaming. 15\,--\,30 Minuten.
  \item \nivLiii{} \textendash{} \fachbegriff{Management.} Decision Architecture, Risk Appetite, Integration in Transformation. 30\,--\,45 Minuten.
\end{itemize}

\noindent Jede Aufgabe enth\"alt eine Niveau-Markierung, einen Zeit-Richtwert, den Aufgabentext, einen Antwort-Freiraum und eine \emph{YouTube-Suchquery}.

\begin{infobox}[Hinweis: Capstone-Tag]
\textbf{Tag 16 ist der Capstone des Kurses.} Die Aufgaben verbinden Inhalte aus den vorhergehenden Tagen: BPMN/EPC (Prozesssicht), RPA-Governance (Audit Trail), Sustainable IT (Datenpipelines), Green IT (Datenqualit\"at), Change Leadership (Verankerung). Das ist beabsichtigt: \emph{ESG-Compliance ist kein Reporting-Projekt, sondern eine integrierte Entscheidungs- und Kontrollarchitektur}.
\end{infobox}

\clearpage

% =========================================================================
% L1 AUFGABEN
% =========================================================================
\section*{Niveau L1 \textendash{} Wissen \& Reproduktion}
\addcontentsline{toc}{section}{L1 -- Wissen}

% --- AUFGABE 1 -----------------------------------------------------------
% LERNPLATT_HILFSVIDEOS
\section*{Hilfsvideos zu diesem Tag}
\addcontentsline{toc}{section}{Hilfsvideos zu diesem Tag}
\thispagestyle{plain}

\noindent Die folgenden YouTube-Erkl\"arvideos vermitteln die Inhalte, die Sie zur Bearbeitung der Aufgaben ben\"otigen. Klicken Sie auf den Titel, um das Video direkt zu \"offnen; die vollst\"andige URL steht in Klammern.

\begin{description}[style=nextline,leftmargin=18pt,labelindent=0pt,itemsep=8pt]
  \item[1.~CSRD, ESRS und EU-Taxonomie — wie hängt das alles zusammen?] \href{https://youtu.be/_ypAyhFA3MQ}{\textcolor{lpAccent}{https://youtu.be/_ypAyhFA3MQ}}\\[2pt]
  {\footnotesize\color{lpGray}(URL: \nolinkurl{https://youtu.be/_ypAyhFA3MQ})}
  \item[2.~EU-Taxonomie kurz erklärt] \href{https://youtu.be/DfoXSnvrVhc}{\textcolor{lpAccent}{https://youtu.be/DfoXSnvrVhc}}\\[2pt]
  {\footnotesize\color{lpGray}(URL: \nolinkurl{https://youtu.be/DfoXSnvrVhc})}
  \item[3.~ESG, CSRD und EU-Richtlinien für Unternehmen] \href{https://youtu.be/vZ-kc1M0oSk}{\textcolor{lpAccent}{https://youtu.be/vZ-kc1M0oSk}}\\[2pt]
  {\footnotesize\color{lpGray}(URL: \nolinkurl{https://youtu.be/vZ-kc1M0oSk})}
  \item[4.~Nachhaltigkeitsberichterstattung in der Praxis — DAX-Beispiel] \href{https://youtu.be/xw55BGza2-g}{\textcolor{lpAccent}{https://youtu.be/xw55BGza2-g}}\\[2pt]
  {\footnotesize\color{lpGray}(URL: \nolinkurl{https://youtu.be/xw55BGza2-g})}
\end{description}
\vspace{0.6em}
\noindent{\footnotesize\color{lpGray}\textit{Tipp:} \"Offnen Sie das Video, lassen Sie es im Hintergrund laufen und arbeiten Sie parallel an der Aufgabe.}

\clearpage

\aufgabenkopf{1}{CSRD \& EU-Taxonomie aus Prozesssicht}{\nivLi}{8\,min}

\noindent \emph{CSRD} (Corporate Sustainability Reporting Directive) und \emph{EU-Taxonomie} verlangen, dass Unternehmen Nachhaltigkeits-Informationen \emph{verl\"asslich}, \emph{vergleichbar} und \emph{auditf\"ahig} berichten. Aus Prozesssicht hei{\ss}t das: Anforderungen erzeugen \emph{Datenbedarfe}, \emph{Nachweise}, \emph{Kontrollen} und \emph{Verantwortlichkeiten}. Wo entstehen Kennzahlen (Energie, Emissionen, CapEx-/OpEx-Anteile, Policies, Lieferkette), wo werden sie validiert, wo freigegeben?

\medskip
\noindent\textbf{Arbeitsauftrag:}

\begin{enumerate}
  \item Erkl\"aren Sie in je 2 S\"atzen, was \emph{CSRD} und \emph{EU-Taxonomie} fordern \textendash{} aus Prozesssicht, nicht aus Jura-Sicht (z.\,B.\ CSRD: doppelte Wesentlichkeit; EU-Taxonomie: Klassifikation wirtschaftlicher T\"atigkeiten nach \"okologischen Kriterien).
  \item Nennen Sie \textbf{drei typische Risiken} aus Prozesssicht: \emph{uneinheitliche Definitionen}, \emph{Dateninseln}, \emph{Sch\"atzungen ohne Transparenz}, \emph{fehlende Verantwortliche}. F\"ur jedes Risiko: ein konkretes Beispiel.
  \item Erkl\"aren Sie in einem Satz, warum CSRD ein \emph{Prozess-} und \emph{kein Dokument-Thema} ist (Dokument-Sicht: ``Wir bauen einen Bericht''; Prozess-Sicht: ``Wir bauen ein System, das den Bericht erzeugt'').
\end{enumerate}

\ytquery{CSRD EU taxonomy process perspective sustainability reporting data governance}

\antwortlinien{14}

\clearpage

% --- AUFGABE 2 -----------------------------------------------------------
\aufgabenkopf{2}{Governance- \& Kontrollsysteme f\"ur ESG-Reporting}{\nivLi}{10\,min}

\noindent F\"ur belastbares ESG-Reporting braucht es drei Schichten: \emph{Governance} (Rollen + Entscheidungswege + Policies + Review-Zyklen), \emph{Controls} (definierte Pr\"ufungen: Plausibilit\"at, Vollst\"andigkeit, Freigaben, Traceability) und \emph{Audit-Readiness} (Nachweisf\"uhrung, Prozessdokumentation, ``Evidence Packs'').

\medskip
\noindent\textbf{Arbeitsauftrag:}

\begin{enumerate}
  \item Skizzieren Sie ein \textbf{Rollenmodell} mit f\"unf Rollen (\emph{ESG Steering}, \emph{Data Owner}, \emph{Control Owner}, \emph{Reporting Factory Lead}, \emph{Internal Audit}). Pro Rolle: \emph{Hauptverantwortung}, \emph{Entscheidungsrechte}, \emph{Eskalationsweg}.
  \item Definieren Sie \textbf{f\"unf Control-Typen} (z.\,B.\ \emph{Completeness Check}, \emph{Plausibilit\"atspr\"ufung mit Schwellen}, \emph{Vier-Augen-Freigabe}, \emph{Change-Log / Versionierung}, \emph{Zugriffskontrolle}). Pro Typ: \emph{wozu da}, \emph{wer pr\"uft}, \emph{wie h\"aufig}.
  \item Definieren Sie \textbf{Audit-Readiness} in einem Satz \textendash{} und nennen Sie \emph{drei Evidence-Pack-Elemente} (z.\,B.\ Datenquelle, Transformationslogik, Freigabeprotokoll, Annahmenregister, Versionierung).
\end{enumerate}

\ytquery{ESG governance controls audit readiness evidence pack CSRD reporting}

\antwortlinien{14}

\clearpage

% --- AUFGABE 3 -----------------------------------------------------------
\aufgabenkopf{3}{KPI-Dictionary \& Datenqualit\"atslabel}{\nivLi}{10\,min}

\noindent Eine \emph{Single Source of Truth} f\"ur ESG-KPIs verlangt ein \emph{KPI-Dictionary}: pro KPI eine \emph{Definition}, \emph{Systemgrenze}, \emph{Berechnung}, \emph{Datenquelle}, \emph{Datenqualit\"at} (A/B/C, siehe Tag 14), \emph{Anti-Gaming-Mechanik}, \emph{Owner}. Materialit\"at-/Risiko-Logik bestimmt, \emph{welche} Controls und KPIs intensiv gef\"uhrt werden \textendash{} ``nicht alles pr\"ufen, sondern das Richtige''.

\medskip
\noindent\textbf{Arbeitsauftrag:}

\begin{enumerate}
  \item Definieren Sie ein \textbf{KPI-Steckbrief-Template} mit mindestens neun Feldern (Name, Definition, Berechnung, Systemgrenze, Einheit, Datenquelle, Datenqualit\"atslabel, Anti-Gaming-Regel, Owner).
  \item F\"ullen Sie das Template f\"ur \textbf{drei KPIs} aus: \emph{Energieverbrauch IT}, \emph{Emissionsintensit\"at IT-Services}, \emph{Anteil verifizierte Vendor-Daten}.
  \item Erkl\"aren Sie in einem Satz, warum ein \emph{KPI ohne Datenqualit\"atslabel nicht im CSRD-Bericht erscheinen darf} \textendash{} und welche Konsequenz das auf die Reporting-Pipeline hat (Pflichtfeld; Pipeline lehnt unmarkierte Werte ab).
\end{enumerate}

\ytquery{KPI dictionary ESG data quality label anti gaming reporting CSRD}

\antwortlinien{14}

\clearpage

% =========================================================================
% L2 AUFGABEN
% =========================================================================
\section*{Niveau L2 \textendash{} Anwendung \& Transfer}
\addcontentsline{toc}{section}{L2 -- Anwendung}

% --- AUFGABE 4 -----------------------------------------------------------
\aufgabenkopf{4}{ESG-Reportingprozess End-to-End in 8\,--\,10 Schritten}{\nivLii}{30\,min}

\begin{infobox}[Ausgangslage]
Ein Unternehmen muss ESG-relevante Daten aus IT, Facility, Einkauf, HR, Finance konsolidieren. Heute existieren Excel-Silos, Definitionen sind unklar.
\end{infobox}

\noindent\textbf{Arbeitsauftrag:}

\begin{enumerate}
  \item Skizzieren Sie einen \textbf{End-to-End-Reportingprozess in 8\,--\,10 Schritten} (textuell): \emph{KPI-Dictionary aktualisieren} \textrightarrow{} \emph{Datenanforderung an Owner} \textrightarrow{} \emph{Datenerfassung (Quellen)} \textrightarrow{} \emph{Validierung (Plausibilit\"at / Completeness)} \textrightarrow{} \emph{Konsolidierung} \textrightarrow{} \emph{Freigabe (Owner + Reporting Lead)} \textrightarrow{} \emph{Ver\"offentlichung / Reporting} \textrightarrow{} \emph{Audit-Trail / Evidence Pack} \textrightarrow{} \emph{Lessons Learned}.
  \item \textbf{Drei Kontrollpunkte} entlang des Prozesses: z.\,B.\ Completeness-Check (vor Konsolidierung); Plausibilit\"atspr\"ufung mit Schwellen (vor Freigabe); Vier-Augen-Freigabe (vor Ver\"offentlichung). Pro Kontrolle: \emph{Owner}, \emph{Eingangskriterium}, \emph{Ausgang/Folge}.
  \item \textbf{F\"unf Evidence-Artefakte}: Datenquelle, Berechnung (Script-Version), Freigabe, Ausnahmebehandlung, Versionierung.
  \item \textbf{Zwei Transparenzma{\ss}nahmen} f\"ur schwache Daten: \emph{Datenqualit\"atslabel} (A/B/C; D nicht ver\"offentlichbar); \emph{Annahmenregister} mit jedem Reporting versendet.
  \item \textbf{Eine Trade-off-Entscheidung:} Welcher Schritt erlaubt eine Sch\"atzung \emph{bewusst} (z.\,B.\ Lieferkettendaten als Label C), und welche Guardrails (Methodik dokumentiert, Verbesserungsplan, Sponsor-Sign-off) sind verpflichtend?
\end{enumerate}

\ytquery{CSRD reporting process end to end controls evidence data quality}

\antwortlinien{24}

\clearpage

% --- AUFGABE 5 -----------------------------------------------------------
\aufgabenkopf{5}{Acht Controls + Evidence Pack Checkliste}{\nivLii}{25\,min}

\begin{infobox}[Fall: IndustriCo Europe (regulatorisch sensibel)]
Mehrere Werke, komplexe Lieferkette; Reporting muss \emph{konsistent}, \emph{nachvollziehbar} und \emph{auditf\"ahig} werden.
\end{infobox}

\noindent\textbf{Arbeitsauftrag:}

\begin{enumerate}
  \item Definieren Sie \textbf{acht Controls} entlang des Reportingprozesses: Pflichtfelder/Completeness; Plausibilit\"atsgrenzen (Schwellen); 2-Augen-Freigabe; \"Anderungslog (wer/was/warum); Zugriffskontrolle; Ausnahmegenehmigung; Reconciliation zu Finance-Daten; Sampling-Review f\"ur Audit-Readiness. Pro Control: \emph{Eingangskriterium}, \emph{Ausl\"oser}, \emph{Owner}, \emph{Eskalation bei Fehler}.
  \item Erstellen Sie eine \textbf{Evidence-Pack-Checkliste} pro KPI: Datenquelle (System/Export), Transformationslogik (ETL), Berechnungsblatt/Script-Version, Freigabe\-protokoll, Ausnahmenregister, Datenqualit\"atslabel, Change-Requests.
  \item Markieren Sie f\"ur jede Control: ist sie \emph{Minimum Viable Control} (MVC) f\"ur den \emph{Dry Run} \textendash{} oder erst \emph{Scale-Phase}? Begr\"undung in je einem Satz.
\end{enumerate}

\ytquery{CSRD controls evidence pack minimum viable controls reconciliation sampling}

\antwortlinien{22}

\clearpage

% --- AUFGABE 6 -----------------------------------------------------------
\aufgabenkopf{6}{KPI-Set 6\,--\,8 + Anti-Gaming-Kopplung}{\nivLii}{25\,min}

\begin{infobox}[Aufgabenstellung]
F\"ur CSRD / EU-Taxonomie soll ein robustes \emph{KPI-Set} (6\,--\,8 KPIs) entwickelt werden. Jeder KPI muss \emph{Datenqualit\"atslabel} und \emph{Anti-Gaming-Mechanik} besitzen. Ohne Kopplung droht Greenwashing (Tag 14).
\end{infobox}

\noindent\textbf{Arbeitsauftrag:}

\begin{enumerate}
  \item Definieren Sie ein \textbf{KPI-Set mit 6\,--\,8 KPIs}, das alle drei ESG-Dimensionen abdeckt (Environment / Social / Governance). Vorschl\"age: \emph{Energieverbrauch} (A/B/C), \emph{Emissionsintensit\"at} (mit Systemgrenze), \emph{Anteil verifizierte Vendor-Daten}, \emph{Datenqualit\"atsindex}, \emph{Control Pass Rate}, \emph{Timeliness (Reporting p\"unktlich)}, \emph{Exception Rate}, \emph{Mitarbeiter-Diversit\"atskennzahl}.
  \item Pro KPI: \emph{Definition / Berechnung / Owner / Ziel / Datenqualit\"at / Anti-Gaming}.
  \item \textbf{Kopplungslogik} f\"ur mindestens drei KPI-Paare: z.\,B.\ ``Kosten- oder CO\textsubscript{2}-KPI z\"ahlt nur, wenn Datenqualit\"atslabel $\geq$\,B und Control Pass Rate $\geq$\,X.'' Erkl\"aren Sie, was diese Kopplung strukturell verhindert.
  \item \textbf{Schutz-KPI:} Eine KPI, die das System \emph{vor sich selbst} sch\"utzt (z.\,B.\ \emph{Anteil Reports mit DQ-Label und Annahmenregister} \emph{$\geq$\,95\,\%}). Owner und Konsequenz bei Unter\-schreiten.
\end{enumerate}

\ytquery{ESG KPI set anti gaming coupling data quality CSRD greenwashing}

\antwortlinien{22}

\clearpage

% --- AUFGABE 7 -----------------------------------------------------------
\aufgabenkopf{7}{Integration in Transformation (RPA, ITSM, Green IT, Change)}{\nivLii}{25\,min}

\begin{infobox}[These (Kurs-Kontext)]
ESG-Compliance darf nicht als isoliertes ``Reporting-Projekt'' laufen. Sie muss in die bereits laufenden Transformationsprogramme dieses Kurses integriert sein: Prozessdigitalisierung (BPMN/EPC, Workflow), RPA-Governance (Audit Trail), Sustainable IT (Service-Portfolio, Capacity), Green IT (Mess-System, Anti-Greenwashing), Change Leadership (Verankerung).
\end{infobox}

\noindent\textbf{Arbeitsauftrag:}

\begin{enumerate}
  \item F\"ur \textbf{jedes der f\"unf Transformationsprogramme} (BPMN/EPC, RPA, ITSM, Green IT, Change): Schreiben Sie je 2\,--\,3 S\"atze, wie dieses Programm \emph{direkt auf CSRD einzahlt} (z.\,B.\ RPA-Audit-Trail $\rightarrow$ liefert nachweisf\"ahige Datenherkunft; Green-IT-Messsystem $\rightarrow$ liefert die Inputs f\"ur Scope-2-Berechnungen).
  \item Identifizieren Sie \textbf{drei konkrete Prozess-/Tooling-Initiativen}, die in den n\"achsten 14 Wochen direkt CSRD-relevant werden (z.\,B.\ DWH/ETL f\"ur ESG-Events; Workflow-Freigaben f\"ur KPI-\"Anderungen; Monitoring \& Dashboards mit Audit Trail).
  \item \textbf{Drei Anti-Muster}, die diese Integration zerst\"oren (z.\,B.\ Reporting wird separat in Excel gebaut; ESG-Team hat keine Anbindung an IT-Architektur; KPI-\"Anderungen ohne Change-Log).
\end{enumerate}

\ytquery{ESG transformation integration RPA ITSM green IT change management CSRD}

\antwortlinien{20}

\clearpage

% --- AUFGABE 8 -----------------------------------------------------------
\aufgabenkopf{8}{Capstone-Fallstudie ``IndustriCo Europe''}{\nivLii}{45\,min}

\begin{infobox}[Capstone-Fallstudie: IndustriCo Europe]
\textbf{Unternehmen:} ``IndustriCo Europe'' \textendash{} mehrere Werke, komplexe Lieferkette.

\smallskip
\textbf{Problem:}
\begin{itemize}
  \item ESG-Daten sind fragmentiert.
  \item Keine einheitlichen Definitionen.
  \item Reporting muss konsistent, nachvollziehbar und auditf\"ahig werden.
  \item Parallel l\"auft ein Transformationsprogramm (Prozessdigitalisierung, Green IT, RPA).
\end{itemize}

\smallskip
\textbf{Restriktionen:} 14 Wochen bis zum internen Dry-Run; Budget 250\,k\,\euro{}; Datenl\"ucken (Energie/Lieferkette); Zielkonflikt Transparenz vs.\ Aufwand; Reputationsrisiko (Greenwashing).
\end{infobox}

\noindent\textbf{Arbeitsauftrag (integriert):}

\begin{enumerate}
  \item \textbf{Reportingprozess} (Plan-Collect-Validate-Consolidate-Approve-Report-Assure) inkl.\ Rollen und Schnittstellen (IT, Einkauf, Werke, Finance).
  \item \textbf{Acht Controls} (siehe Aufgabe 5) mit Owner und Eskalation.
  \item \textbf{Evidence Pack Checkliste} pro KPI (siehe Aufgabe 5).
  \item \textbf{KPI-Set (6\,--\,8 KPIs)} inkl.\ Datenqualit\"atslabel und Anti-Gaming-Kopplung (siehe Aufgabe 6).
  \item \textbf{Integration in Transformation}: Drei Prozess-/Tooling-Initiativen, die direkt einzahlen (DWH/ETL, Workflow-Freigaben, BI-Dashboards mit Audit Trail; RPA f\"ur Datensammlung mit Governance).
  \item \textbf{Entscheidung unter Unsicherheit:} Welche \emph{zwei Datenl\"ucken} akzeptieren Sie im Dry Run \textendash{} und welche Kontrollen / Transparenzma{\ss}nahmen kompensieren das? Pro Entscheidung: Annahme, verworfene Alternative, Guardrails, Verbesserungsplan, Fr\"uhwarnsignal.
\end{enumerate}

\ytquery{CSRD case study reporting process controls KPI transformation industrial}

\antwortlinien{36}

\clearpage

% =========================================================================
% L3 AUFGABEN
% =========================================================================
\section*{Niveau L3 \textendash{} Management \& Entscheidungsarchitektur}
\addcontentsline{toc}{section}{L3 -- Management}

% --- AUFGABE 9 -----------------------------------------------------------
\aufgabenkopf{9}{Risk Appetite \& Vorstands-Transparenzregeln}{\nivLiii}{35\,min}

\begin{infobox}[Rolle: CSO / CFO / CIO gemeinsam]
\textbf{Ihre Rolle:} \emph{Chief Sustainability Officer} im Tandem mit \emph{CFO} und \emph{CIO}; Sie verantworten gemeinsam die ESG-Governance.

\smallskip
\textbf{Auftrag:} In 16 Wochen eine auditf\"ahige ESG-Governance etablieren, die CSRD und EU-Taxonomie in Prozesse, Daten, Controls, KPIs und Tooling \"ubersetzt \textendash{} und dauerhaft in Transformationsprogramme integriert ist.
\end{infobox}

\noindent\textbf{Arbeitsauftrag:}

\begin{enumerate}
  \item \textbf{Risk Appetite f\"ur ESG-Reporting:} F\"unf Bereiche, in denen Sie \emph{bewusst} Risiko akzeptieren (z.\,B.\ Lieferkettendaten als DQ-C, Sch\"atzungen f\"ur Scope-3, Sampling statt 100-\%-Pr\"ufung). Pro Bereich: \emph{Annahme}, \emph{Risiko-Quantifizierung} (qualitativ + grob), \emph{Guardrail}, \emph{Fr\"uhwarnsignal}, \emph{Eskalation}.
  \item \textbf{Drei Vorstands-Transparenzregeln:} z.\,B.\ ``Jede ver\"offentlichte Zahl tr\"agt das schlechteste Datenqualit\"atslabel im Aggregat''; ``Annahmenregister geht im selben Versand mit''; ``D-Werte sind nicht ver\"offentlichbar''.
  \item \textbf{Drei nicht delegierbare Entscheidungen} f\"ur den Vorstand: z.\,B.\ \emph{Systemgrenzen}, \emph{Risk Appetite f\"ur Sch\"atzungen}, \emph{Transparenzregeln gegen Greenwashing}. Erkl\"aren Sie pro Entscheidung, warum Delegation strukturell fehlerhaft w\"are.
  \item \textbf{Eine Vorstandsvorlage (1 Absatz):} Wie pr\"asentieren Sie das Risk Appetite Framework im n\"achsten Vorstandstreffen \textendash{} mit klarem Anker, was Sie entscheiden lassen, und was Sie als Empfehlung tragen?
\end{enumerate}

\ytquery{ESG risk appetite governance board transparency rules CSRD audit committee}

\antwortlinien{26}

\clearpage

% --- AUFGABE 10 ----------------------------------------------------------
\aufgabenkopf{10}{Minimum Viable Controls (MVC) f\"ur Dry Run}{\nivLiii}{30\,min}

\noindent\textbf{Diskussionsaufgabe.}

\noindent Im Dry Run (4\,--\,6 Wochen) kann nicht \emph{jede} Control implementiert sein \textendash{} aber das Ergebnis muss \emph{trotzdem} belastbar sein. Die Senior-Frage: \emph{Welche Controls sind ``Minimum Viable Controls'' (MVC), welche sind Overkill}?

\medskip
\noindent\textbf{Arbeitsauftrag:}

\begin{enumerate}
  \item Definieren Sie \textbf{f\"unf MVC} (\emph{Minimum Viable Controls}) f\"ur den Dry Run \textendash{} ohne die das Reporting nicht aussagekr\"aftig ist (z.\,B.\ Completeness-Check, 2-Augen-Freigabe f\"ur Top-3-KPI, DQ-Labeling-Pflicht, Annahmenregister, Change-Log Pflicht).
  \item Definieren Sie \textbf{drei Controls}, die in den \emph{Scale}-Phasen (Wo 7\,--\,12) dazukommen, im Dry Run aber bewusst ausgesetzt sind (z.\,B.\ Reconciliation zu Finance-Daten f\"ur alle KPIs; Sampling-Review f\"ur Vendor-Daten; Auto-Validierung gegen historische Werte).
  \item \textbf{Kompensationsmechanik:} Welche \emph{f\"unf Ma{\ss}nahmen} f\"angt das Fehlen der drei Scale-Controls im Dry Run ab (z.\,B.\ \emph{Sampling per Hand statt automatisiert}; \emph{Sponsor-Sign-off statt Reconciliation}; \emph{Audit-Trockenlauf in Wo 10})?
  \item \textbf{Stoppkriterien:} Drei Schwellen, die einen Dry Run als ``nicht ausreichend'' kennzeichnen (z.\,B.\ $>$\,3 KPIs ohne DQ-Label; Control Pass Rate $<$\,80\,\%; mehr als ein materieller Audit-Finding bei Stichprobe).
  \item \textbf{Faustregel} (1 Satz): Wann ist eine Control \emph{MVC}, wann \emph{Overkill}?
\end{enumerate}

\ytquery{minimum viable controls dry run sampling audit readiness CSRD scale phase}

\antwortlinien{20}

\clearpage

% --- AUFGABE 11 ----------------------------------------------------------
\aufgabenkopf{11}{Capstone-Transferprojekt \textendash{} CSRD/EU-Taxonomie Decision Architecture}{\nivLiii}{60\,min}

\begin{infobox}[Rolle und Ausgangslage]
\textbf{Ihre Rolle:} \emph{Chief Sustainability Officer} im Tandem mit \emph{CFO} und \emph{CIO} (gemeinsame Steuerung).

\smallskip
\textbf{Auftrag:} In 16 Wochen eine auditf\"ahige ESG-Governance etablieren, die EU-Taxonomie- und CSRD-Anforderungen in Prozesse, Daten, Controls, KPIs und Tooling \"ubersetzt \textendash{} und in laufende Transformationsprogramme (BPMN/EPC, RPA, Sustainable IT, Green IT, Change) integriert ist.

\smallskip
\textbf{Restriktionen:} Datenl\"ucken; heterogene Systeme; begrenzte Kapazit\"at; hoher Zeitdruck; reputations\-kritisches Umfeld; Zielkonflikte (Transparenz vs.\ Aufwand, Speed vs.\ Assurance).
\end{infobox}

\noindent\textbf{Capstone-Arbeitsauftrag:} Entwickeln Sie eine vollst\"andige \emph{Entscheidungsarchitektur} in 9 Schritten. Pr\"asentationstauglich, kein Fl\"ie{\ss}text-Roman.

\begin{enumerate}
  \item \textbf{Top-12 Entscheidungen:} Systemgrenzen; KPI-Dictionary Ownership; Datenqualit\"atslabels; Control-Set; Ausnahmeprozess; Tooling-Architektur (DWH/ETL/BI); Reporting-Cadence; Audit-Readiness; Vendor-Datenstrategie; Kommunikations-/Transparenzregeln; Priorisierung; Budget/Ressourcen.
  \item \textbf{Governance \& Operating Model:} ESG Steering, Data Owners, Control Owners, Reporting Factory, RACI, Eskalation, Entscheidungsprotokolle.
  \item \textbf{Controls \& Evidence Framework:} Minimum Viable Controls $\rightarrow$ Scale; Evidence Packs; Change-/Versioning; Access Control; Sampling-/Assurance-Plan.
  \item \textbf{Integration in Transformation:} Verankerung in Prozessdigitalisierung (Workflows), Process Intelligence (Monitoring), Green IT Ma{\ss}nahmen, Vendor Management; KPI-Kopplungen gegen Greenwashing.
  \item \textbf{Roadmap:} Dry Run (4\,--\,6 Wochen) \textrightarrow{} Stabilisierung (6\,--\,10 Wochen) \textrightarrow{} Assurance Readiness (laufend), inkl.\ Change Management \& Training.
  \item \textbf{Risikomodell:} Top-5 Risiken (Greenwashing, Datenqualit\"at, fehlende Verantwortlichkeit, Audit-Findings, KPI-Gaming) + Gegenma{\ss}nahmen.
  \item \textbf{Zwei Entscheidungen unter Unsicherheit:}
    \begin{itemize}
      \item Welche KPI / Reports werden trotz Datenqualit\"at C ver\"offentlicht (Transparenzentscheidung) \textendash{} und welche Guardrails sind Pflicht?
      \item Welche Controls werden im Dry Run reduziert \textendash{} und wie wird das Risiko kompensiert (Sampling, Zusatzreview, Board-Freigabe)?
    \end{itemize}
    Pro Entscheidung: Annahme, verworfene Alternative, Fr\"uhwarnsignal, Reviewdatum.
  \item \textbf{Anschluss an Strategie:} F\"unf Argumente, warum diese Entscheidungsarchitektur Strategie tr\"agt (Compliance, Reputation, Investorenkommunikation, Effizienz, Anschluss an Transformation).
  \item \textbf{Pr\"asentationssatz f\"ur Vorstand und Aufsichtsrat} (2\,--\,3 S\"atze): Was ist die Kernbotschaft, die das Mandat tr\"agt, ohne in Details abzugleiten?
\end{enumerate}

\noindent\textbf{Bewertungskriterien (Senior-Capstone):} Anschluss an Strategie \textperiodcentered{} Integration mit Tagen 1\,--\,15 \textperiodcentered{} Pragmatismus unter Restriktionen \textperiodcentered{} Reife des Entscheidens unter Unsicherheit \textperiodcentered{} Anti-Greenwashing-Logik \textperiodcentered{} Auditf\"ahigkeit \textperiodcentered{} Wirkung: \emph{nicht ``Reporting erf\"ullen'', sondern ein verantwortetes Governance-, Daten- und Entscheidungs-Betriebssystem aufbauen}.

\ytquery{CSRD EU taxonomy decision architecture CSO CFO CIO governance integration}

\antwortlinien{44}

\clearpage

\section*{Abschluss \& Kursabschluss}
\thispagestyle{plain}

\noindent Sie haben alle elf Aufgaben des Capstone-Tags 16 bearbeitet \textendash{} und damit den Kurs \emph{Digital Business Modelling \& Transformation} (Kurs 71 / Dig 42) durchlaufen. Vergleichen Sie nun Ihre Antworten mit dem \emph{L\"osungsheft Tag 16}.

\medskip
\begin{infobox}[Was Sie jetzt k\"onnen sollten]
\begin{itemize}
  \item CSRD und EU-Taxonomie aus Prozesssicht denken \textendash{} nicht als Dokument-Projekt.
  \item Eine Governance-Schicht aus Rollen, Controls und Audit-Readiness aufsetzen.
  \item Ein KPI-Dictionary mit Datenqualit\"atslabels und Anti-Gaming-Kopplung bauen.
  \item Einen End-to-End-Reportingprozess mit Controls und Evidence Pack designen.
  \item Minimum Viable Controls von Overkill unterscheiden \textendash{} im Dry Run wie in Scale.
  \item ESG-Compliance in laufende Transformationsprogramme (BPMN/EPC, RPA, ITSM, Green IT, Change) integrieren, statt isoliert zu reportieren.
  \item Vorstands-Transparenzregeln gegen Greenwashing formulieren und durchsetzen.
  \item Eine vollst\"andige Capstone-Entscheidungsarchitektur tragen, die unter Unsicherheit handlungsf\"ahig bleibt.
\end{itemize}
\end{infobox}

\medskip
\noindent\textbf{Schlussgedanke (Kurs 71).} Der Bogen dieses Kurses war eine Bewegung von \emph{Notation} (BPMN, EPC) \"uber \emph{Automatisierung} (RPA, Governance) und \emph{Steuerung} (ITSM, Green IT) zu \emph{F\"uhrung} (Change Leadership) und \emph{Verantwortung} (CSRD). Diese f\"unf Ebenen sind \emph{eine} Architektur, kein Werkzeugkasten. Wer Tag 16 sauber tr\"agt, hat Tag 1 bis 15 verstanden \textendash{} und umgekehrt.

\vspace{1cm}
\noindent{\color{lpAccent}\rule{\linewidth}{0.6pt}}\\[4pt]
{\footnotesize\sffamily\color{lpGray}Lernplatt \textperiodcentered{} by Can Siebert \textperiodcentered{} Kurs 71 (Dig 42) \textperiodcentered{} Tag 16 \textperiodcentered{} Aufgabenheft (Capstone) \textperiodcentered{} \textcopyright{} 2026}

\end{document}
